%==============================================================================
%  DATOS DEL TRABAJO
%
%  - El título tiene que coincidir con el del TFG asignado en Ebrón.
%  - Estos datos no salen en el PDF: la portada la genera Ebrón al solicitar
%    la defensa (con lo que tenga registrado). Se dejan por si algún día
%    quieres la portada de la clase para imprimir.
%==============================================================================

\title{Adaptación del modelo OWSM para la decodificación de actividad neuronal intracortical en el reto Brain-to-Text '25}
\author{Kiril Ivaylov Tzenkov}
\tutor{David Gimeno Gómez\\Salvador España Boquera}
\curs{2026-2027}

% Palabras clave: entre 3 y 10 (guía de la ETSINF para Ciencia de Datos).
% Es lo único que usa la clase para las palabras clave: \keywords con tres
% argumentos, en este orden: valenciano, castellano e inglés. Cada resumen
% saca las de su idioma.
\keywords{interfície cervell-ordinador, descodificació de la parla, esclerosi lateral amiotròfica, activitat neuronal intracortical, reconeixement automàtic de la parla, aprenentatge per transferència, variabilitat entre sessions, OWSM, Brain-to-Text '25}
         {interfaz cerebro-computador, decodificación del habla, esclerosis lateral amiotrófica, actividad neuronal intracortical, reconocimiento automático del habla, aprendizaje por transferencia, variabilidad entre sesiones, OWSM, Brain-to-Text '25}
         {brain--computer interface, speech decoding, amyotrophic lateral sclerosis, intracortical neural activity, automatic speech recognition, transfer learning, cross-session variability, OWSM, Brain-to-Text '25}
